% !TEX root = ../Main.tex
\chapter{电源内阻}

真实的电源不是"理想"的：它自己也有一点内阻 $r$。电源标称的电压 $U$（电动势）并不会全部加到外电路上，有一部分被内阻"吃掉"了。

\section{路端电压 \texorpdfstring{$U'=U-Ir$}{U'=U-Ir}}

把电源看成"理想电压 $U$ $+$ 串联内阻 $r$"。当有电流 $I$ 流过时，内阻上分掉 $Ir$，真正送到两端（外电路）的电压只剩

\[
U'=U-Ir.
\]

\begin{center}
\begin{circuitikz}[scale=0.9, american]
    \draw (0,0) to[battery1, l=$U$] (0,2.2);
    \draw (0,2.2) to[R, l=$r$] (2.2,2.2);
    \draw (2.2,2.2) -- (3.0,2.2);
    \draw (0,0) -- (3.0,0);
    \draw[->] (3.0,2.2) -- (3.0,0) node[midway, right] {$U'$};
\end{circuitikz}
\end{center}

这就是\textbf{路端电压}：它随电流变化，电流越大、内阻分压越多、路端电压越低。

\section{接上负载求电流}

给电源接一个负载电阻 $R_L$（图中的灯），用电压表测路端电压。这时电压表读到的 $U_1$ 等于负载两端电压 $U_L$，也就是路端电压：

\[
U_1=U_L=U-Ir.
\]

\begin{center}
\begin{circuitikz}[scale=0.9, american]
    \draw (0,0) to[battery1, l={$U,\ r$}] (0,2.2);
    \draw (0,2.2) -- (3.2,2.2);
    \draw (3.2,2.2) to[lamp, l=$R_L$] (3.2,0);
    \draw (3.2,0) -- (0,0);
    \draw (1.4,2.2) to[voltmeter, l=$V_1$] (1.4,0);
\end{circuitikz}
\end{center}

整个回路是 $U$、$r$、$R_L$ 串联，由欧姆定律得电流

\[
I=\frac{U}{r+R_L}.
\]

\section{变阻器：滑动时谁升谁降}

把负载换成滑动变阻器 $R$，再串一个电流表。\textbf{减小 $R$}（$R\downarrow$）时，各量怎么变？

\begin{center}
\begin{circuitikz}[scale=1.0, american]
    \draw (0,0) to[battery1, l={$U,\ r$}] (0,2.4);
    \draw (0,2.4) -- (2.2,2.4) to[ammeter, l=$A_1$] (4.4,2.4);
    \draw (4.4,2.4) -- (4.4,0);
    \draw (0,0) -- (4.4,0);
    \draw (1.4,2.4) to[voltmeter, l_=$V_1$] (1.4,0.9);
    \draw (1.4,0.9) -- (1.4,0.6) to[vR, l=$R$] (3.0,0.6) -- (3.0,0);
\end{circuitikz}
\end{center}

由 $I=\dfrac{U}{r+R}$：$R$ 减小，分母变小，电流
\[
I=\frac{U}{r+R}\ \uparrow.
\]
再由 $U_1=U-Ir$：电流增大，内阻分压 $Ir$ 增大，路端电压
\[
U_1=U-Ir\ \downarrow.
\]

\rmkb{
    一根主线贯穿全章：$U'=U-Ir$。$R\downarrow\Rightarrow I\uparrow\Rightarrow Ir\uparrow\Rightarrow$ 路端电压 $U_1\downarrow$。内阻 $r$ 越大，这种"越用劲、电压掉得越狠"的现象越明显。
}
